\section{Discussion and Limitations}
The full-test comparison establishes that feature sharing benefits the
evaluated pretrained student under the stated conditions. It holds model
weights fixed and changes collaboration availability. It does not determine
the ranking of separately trained local, early-fusion, and intermediate-fusion
models. In particular, the planned early-fusion baseline remains unmeasured.

The scaling experiments address a different question: whether complete
collaboration is computationally expensive as the number of slots increases,
and what happens when its fine interactions are restricted. The tree ablation
shows a measured runtime--accuracy trade-off. Removing fine edges can reduce
access to complementary observations in a single fusion round. This is a
plausible explanation of the AP loss, consistent with the mechanism, rather
than a measured causal decomposition of every missed detection.

Several limits constrain these conclusions:
\begin{itemize}
  \setlength{\itemsep}{0pt}
  \setlength{\parsep}{0pt}
  \item \textbf{Ideal communication and poses.} The experiments do not vary
    latency, packet loss, bandwidth limits, asynchronous sensing, or pose
    errors. Their forward times cannot be used directly as wireless latency.
  \item \textbf{Replay beyond five inputs.} Larger workloads increase active
    model slots without adding independent viewpoints or real vehicles.
  \item \textbf{Different evaluation populations.} The full-test AP pools
    4,000 active views; the tree AP uses 30 fixed ego views. Their numerical
    levels are not interchangeable estimates of the same population.
  \item \textbf{One checkpoint and device.} No retraining, model-family
    comparison, multi-GPU experiment, or cross-hardware generalization is
    established. The sparse graph was not trained to compensate for missing
    fine links.
  \item \textbf{Limited uncertainty analysis.} Repeated timings characterize
    runtime variation on a fixed cohort. They do not create new independent
    detection samples, and no confidence interval for population AP is
    reported.
  \item \textbf{Version and scope.} The evaluation uses V2X-Sim 2.0, unlike
    the original paper's 1.0 benchmark. Fixed-checkpoint losses are not a
    training history or evidence about convergence.
\end{itemize}

Follow-up experiments could hold the evaluation population and postprocessing
fixed while varying retained neighbors, selection resolution, and fusion
rounds. Further recordings would allow study beyond five independent
vehicles. Fine-tuning a sparse student and testing communication impairments
would require new training and an explicit transport model, respectively.

\section{Conclusion}
DiscoNet combines per-vehicle feature extraction, pose-based alignment,
spatial attention, and decoder skip connections to produce cooperative car
detections. On the complete prepared V2X-Sim 2.0 test split, enabling sharing
in the same student raises AP by \ReportGainFifty\ and \ReportGainSeventy\
percentage points at IoU 0.5 and 0.7. This directly answers the completed
own-scan/cooperative comparison under ideal conditions.

The separate workload sweep shows increasing full-group computation as
model input count grows. The matched tree experiment reduces fine pair
evaluations but also reduces detection quality: at 15 slots, median runtime
is \ReportTreeReduction\% lower, with AP losses of \ReportTreeDropFifty\ and
\ReportTreeDropSeventy\ points. The results support studying collaboration
sparsity as a trade-off, while preserving the distinction between additional
computation and additional independently observed information.
